% page 154 du fac-simile (image p-153 du scan)
% bloc 154.9 x 233.3 mm ; marge gauche 8.0 mm
\pgc{-12.700mm}{0.000mm}{
\l{144}
\l{}
\l{}
\l{\sou{eskorto}.\cel{2}Trupo qua akompanas persono de loko ad altra,}
\l{\cel{3}por surveyar lua sekureso o por honorizar lu per lo;}
\l{\cel{3}kaptiton, por impedar lua eskapo ; konvoyon, por pro-}
\l{\cel{3}tektar lu. - DEFIRS.}
\l{}
\l{\sou{eskoto}. (\sou{navo}). Grosa kordego manuo-movebla, fixigita an}
\l{\cel{3}la angulo di basa segli por desvolvar od extensar li.}
\l{\cel{3}- DFS.}
\l{}
\l{\sou{eskrokar}. (\sou{trans.}) Furtar (ulo de ulu) per dupigar (lu).}
\l{\cel{3}FI.}
\l{}
\l{\sou{eskudelo}. Vazo kava, ligna, terakota o metala, e c., en qua}
\l{\cel{3}onu servas buliono, supo, por manjar li. -FIS.}
\l{}
\l{\sou{eskupar}. (\sou{trans.}) Vakuigar per ligna shovelo kava (la aquo}
\l{\cel{3}qua eniras la navi). - EF.}
\l{}
\l{\sou{-esm-}.\cel{2}Sufixo qua signifikas la numero.}
\l{}
\l{\sou{esmerilio}. (\sou{zool.}) Speco di falkono, remarkinda pro la mi-}
\l{\cel{3}kreso di lua staturo, la lejereso di lua flugo e la}
\l{\cel{3}rapideso di lua movi. - L.\cel{1}\sou{falco subbuteo}. - FIS.}
\l{}
\l{\sou{esotera}. (\sou{filoz.}) Qua esas la objekto di docado partikulara,}
\l{\cel{3}intima. - DEFIRS.}
\l{}
\l{\sou{espado}. Atak-armo, ek lamo akuta ek stalo, muntita aden man-}
\l{\cel{3}che konkoza, e quan onu +weras an la flanko, en gaino.}
\l{\cel{3}- FIS.}
\l{}
\l{\sou{espalero}. Muro garnisita maxim-multa-kaze ye lato-greto}
\l{\cel{3}alonge qua onu plantacas frukt-arbori di qui la branchi}
\l{\cel{3}esas aplikita a lu. - DEFRS.}
\l{}
\l{\sou{esparo}. (\sou{navo)}. Stango, min forta kam la masteto, qua su-}
\l{\cel{3}portas la paviliono, ek qua onu facas la stangi per qui}
\l{\cel{3}onu impedas la altra navi venar kolizionar od abordar}
\l{\cel{3}atakoze. - DEF.}
\l{}
\l{\sou{esparseto}. (\sou{bot.}) Planto foreja, analoga a la luzerno, di}
\l{\cel{3}qua la floro esas granda, purpuratra e flavatra, e dis-}
\l{\cel{3}pozita spike od axo-grape. - DEF. L.\cel{1}\sou{onobrychis sativa}.}
\l{}
\l{\sou{esparto}. (\sou{bot.}) Gramineo qua kreskas en la nordo-parto di}
\l{\cel{3}Afrika, di qua la stipeti uzesas por facar objekti diversa}
\l{\cel{3}(chapelo-frami, e c.), papero, e c. L.\cel{1}\sou{stipa tenacissima}.}
\l{\cel{3}DEFIS.}
\l{}
\l{\sou{espelar}. (\sou{trans.}) Lektar la vorti literope, e formacante la}
\l{\cel{3}silabi. - EF.}
}
